\section{The \texttt{Mirror\_toroid} McXtrace Component}
Toroidal shape mirror (in XZ)

\subsection*{Identification}
\begin{itemize}
  \item \textbf{Author:} Erik B Knudsen
  \item \textbf{Origin:} DTU Physics
  \item \textbf{Date:} Jul 2016
\end{itemize}

\subsection*{Description}
This is an implementation of a toroidal mirror which may be curved in two dimensions. To avoid solving quartic equations, the intersection is compited as a combination of two intersections. First, the ray is intersected with a cylinder to catch (almost) the small radius curvature. Secondly, the ray is the intersected with an ellipsoid, with the curvatures matching that of the torus.

In the first incarnation we assume the mirror to be curving outwards (a bump).

Example: Mirror\_toroid(zdepth=0.340,xwidth=0.020,radius=246.9254,radius\_o=246.9254,R0=1)

\subsection*{Input parameters}
Parameters in \textbf{boldface} are required; the others are optional.

\begin{longtable}{p{0.22\textwidth}p{0.12\textwidth}p{0.46\textwidth}p{0.14\textwidth}}
\toprule
\textbf{Name} & \textbf{Unit} & \textbf{Description} & \textbf{Default} \\
\midrule
\endhead
coating & str & Datafile containing either mirror material constants or reflectivity  numbers. & "" \\
zdepth & m & Length of mirror. & 0.1 \\
xwidth & m & Width of mirror. & 0.01 \\
\textbf{radius} & m & Curvature radius &  \\
\textbf{radius\_o} & m & Curvature radius, outwards &  \\
R0 & 1 & Reflectivity of mirror. & 0 \\
\bottomrule
\end{longtable}

\subsection*{Links}
\begin{itemize}
  \item Component source code found in file \texttt{Mirror\_toroid.comp}.
\end{itemize}
\IfFileExists{optics/Mirror_toroid_static.tex}{\input{optics/Mirror_toroid_static.tex}}{}